% !TEX root = ../main.tex

{
\newfontface{\arial}{Arial}[Scale=0.94]
\ctexset{chapter/format+={\arial}}
\chapter{Optimization and Deployment of Instant Communication Streaming Media Congestion Control Algorithm Based on Offline Reinforcement Learning}
\addcontentsline{toc}{chapter}{Optimization and Deployment of Instant Communication Streaming Media Congestion Control Algorithm Based on Offline Reinforcement Learning}

The performance of congestion control algorithms in Real-Time Communication (RTC) directly determines the quality of service and user experience in video calls. With the rapid development of cloud office, remote teaching, online healthcare, cloud gaming, and immersive interactive applications, RTC has evolved from a simple video calling tool into a critical component of digital infrastructure. In such applications, user quality requirements exhibit distinctive multi-objective characteristics: on the one hand, the highest possible media bitrate is needed to ensure video clarity; on the other hand, end-to-end latency, jitter, and freezing rate must be maintained at low levels to ensure real-time interactivity and continuity. Therefore, whether a congestion control algorithm can timely and stably adjust the sending rate in complex network environments directly determines the service quality and user experience of RTC systems.

The Google Congestion Control (GCC) algorithm, widely adopted in current WebRTC systems, essentially belongs to heuristic rule-based congestion control methods. Such algorithms rely on preset thresholds, delay gradients, and packet loss feedback for bitrate adjustment, and exhibit good engineering usability in relatively stable network environments. However, as wireless networks become the primary access scenario for RTC, the time-varying, random, and heterogeneous nature of network links continues to intensify, exposing the insufficient adaptability of heuristic algorithms. In environments such as Wi-Fi and cellular networks, bandwidth may fluctuate drastically within short periods, and links may also experience non-congestive packet loss, sudden jitter spikes, and transient measurement anomalies. Congestion control strategies driven by fixed rules often struggle to distinguish between "real congestion" and "wireless disturbances" in a timely manner, leading to delayed rate adjustment, excessive oscillation, or underutilization of bandwidth, ultimately affecting the smoothness and stability of media streams.

In recent years, reinforcement learning (RL) has provided a new research perspective for congestion control. Compared to traditional methods, RL can learn complex nonlinear decision policies from historical states and feedback, and can simultaneously consider multiple objectives such as throughput, latency, packet loss, and smoothness through reward functions, which is closer to the optimization requirements of real user experience. Online reinforcement learning methods represented by OnRL have demonstrated the feasibility of learning bitrate control policies in mobile video calls. However, online RL typically relies on continuous interaction with the real environment to complete exploration and updates, which means the training process may directly interfere with ongoing RTC sessions, bringing significant trial-and-error costs and deployment risks. Additionally, network environments exhibit obvious non-stationary characteristics, and exploration samples in real systems struggle to systematically cover long-tail scenarios, limiting the generalization ability of policies.

Offline reinforcement learning (Offline RL) provides a feasible new path for the above problems. Its core idea is to complete policy learning directly using historical trajectory data without continuous interaction with the environment. For RTC scenarios, a large number of network states, control actions, and feedback results generated by behavioral policies such as GCC themselves are high-value offline training resources. Recent studies such as Mowgli, Tan et al., and Du et al. have preliminarily verified that passive learning or offline training using historical data can effectively improve the accuracy of bandwidth prediction and the robustness of bitrate control. Through cleaning, modeling, and policy optimization of these trajectories, offline RL is expected to achieve performance improvements over traditional heuristic congestion control strategies without affecting online user experience.

However, existing research still has two prominent problems. First, many reinforcement learning methods rely on online interactive training, where the training process requires constant exploration of new action policies, and their exploratory behavior may significantly degrade real call quality, making them difficult to directly apply in production environments. Second, much existing work mainly relies on network simulation platforms (such as ns-3, SparkRTC) for training and validation. Although simulation tools facilitate variable control and large-scale experiments, there is still an obvious Sim-to-Real gap between simulation environments and real WebRTC systems in terms of wireless channel disturbances, terminal processing delays, system scheduling, and browser implementation details. The performance achieved by policies in simulations may not stably transfer to real devices and real links. The community has launched open-source experimental platforms such as OpenNetLab and AlphaRTC to attempt to bridge this gap, but the public network test nodes have ceased maintenance, making it difficult to continue supporting real network experiments.

Against this background, this thesis addresses the problems of insufficient adaptability, high online learning risks, and limited credibility of simulation validation in existing RTC congestion control methods in real dynamic network environments. Taking real WebRTC communication scenarios as the research object, this thesis constructs an offline reinforcement learning-based congestion control optimization and deployment verification scheme. The research work is carried out around five key components: platform construction, dataset building, model training, inference deployment, and effect evaluation, forming a complete end-to-end technical pipeline.

First, this thesis constructs a WebRTC data collection and validation platform. Based on browser-side WebRTC capabilities (referencing open-source implementations of Hairpin and AlphaRTC), this platform is built for dataset collection and A/B comparative validation, implementing media stream establishment, periodic network state sampling, trajectory preservation, and scenario switching control. The platform supports both basic data collection links and online A/B test links, enabling video call testing between two physical devices, and using NetEm to superimpose controllable bandwidth limits and delay jitter on physical links, recording key indicators such as RTT, jitter, packet loss rate, received throughput, and GCC estimated bandwidth, providing a source of real samples for offline training. The platform's overall design philosophy is based on three core principles: "real physical channel bearing", "end-to-end data closed loop", and "strict physical isolation between collection and evaluation".

Second, this thesis constructs an offline reinforcement learning training dataset. The collected WebRTC trajectories are transformed into state-action-reward-next-state samples required for reinforcement learning. Using a fixed-length historical window as the state representation, features including sending rate, receiving rate, RTT, packet loss rate, and jitter are selected to form the input vector, with the current estimated bandwidth as the action, and a QoE reward function that comprehensively considers throughput gain, latency penalty, packet loss penalty, and action smoothness is designed. Aiming at the problems of invalid zero points, short gaps, and long breakpoints that may exist in the original data, cleaning and interpolation mechanisms are introduced in the data preprocessing stage to improve the quality of training samples and the stability of state representation. The final dataset contains 109,393 valid transition samples, with a state dimension of 50, covering 1,191 continuous valid trajectory segments across multiple real network scenarios.

Third, this thesis uses the Implicit Q-Learning (IQL) algorithm to complete offline policy training. Based on the constructed offline dataset, this thesis selects the IQL algorithm to train the congestion control policy network. This method does not require explicit querying of the value of out-of-distribution actions during policy updates, thus mitigating the distribution shift problem in offline reinforcement learning to a certain extent. After training, the Actor model and corresponding state normalization configuration are exported for subsequent online inference deployment. The neural network architecture adopts a state encoder based on Gated Recurrent Unit (GRU), including a linear encoding layer, a two-layer GRU sequence modeling module, intermediate fully connected layers, and two residual blocks to enhance gradient flow in deep networks.

Fourth, this thesis designs and implements an online inference deployment scheme. To balance the convenience of joint debugging and real-time control performance, this thesis implements two inference methods: one is a remote inference interface based on Python services, supporting HTTP and WebSocket calls; the other is exporting the trained Actor model to ONNX format and completing local inference on the browser side. The former facilitates rapid verification of model input-output logic, while the latter avoids additional latency caused by remote requests and is more suitable for directly driving RTCRtpSender.setParameters() to complete maxBitrate adjustment in the real-time control loop. However, during practical deployment, a fundamental engineering contradiction is quickly exposed when applying the remote inference scheme to real A/B test scenarios: the TCP/WebSocket link required for inference requests shares the same physical channel as the WebRTC media stream being tested. When the channel enters a congested state, the inference requests themselves experience significant queuing delays or even packet loss retransmissions, causing policy decisions to lag by tens or even hundreds of milliseconds precisely when rapid responses are needed most. This finding drives this thesis to adopt browser-local ONNX edge inference as the primary deployment modality for A/B testing.

Fifth, this thesis establishes a QoE-based A/B test evaluation process. Taking GCC as the baseline strategy and the offline RL strategy as the experimental strategy, comparative experiments are conducted on the same WebRTC platform and network scenarios, and algorithm performance is evaluated from multiple dimensions including average received bandwidth, RTT high-percentile values, packet loss rate, jitter, bitrate change smoothness, and comprehensive QoE scores, verifying the actual benefits of the proposed scheme in real application links. A/B testing is conducted across nine representative network scenarios: baseline, 3G, LTE, DSL, low-bandwidth, high-loss, high-delay, fluctuating, and very-bad conditions.

The experimental results show that the offline RL strategy significantly reduces bitrate action oscillation (42\%--73\%) across all scenarios, effectively lowers RTT tail latency (10\%--16\%) and packet loss rate (18\%--25\%) in weak-network scenarios including 3G, LTE, low-bandwidth, and high-delay conditions, achieving positive overall QoE gains. Meanwhile, the policy exhibits certain bandwidth underutilization in high-quality network scenarios, suggesting the need for further optimization in reward weighting and data coverage. In the LTE scenario, the RL strategy reduces RTT P95 from 1410ms to 1258ms while reducing cap\_delta by 73\%, improving comprehensive QoE from -1.115 to -0.926, a difference of +0.189, with QoE improvement approaching statistical significance.

The main contributions of this thesis are as follows. First, constructing a complete end-to-end closed loop from real data collection to online deployment validation, verifying the feasibility of offline reinforcement learning for optimizing congestion control in real WebRTC systems. Second, revealing the importance of the coupling relationship between inference latency and the controlled system in real-time control systems: when inference requests share the physical channel with the media stream under test, the remote inference scheme has fundamental deficiencies in bandwidth-constrained scenarios, an engineering finding that provides enlightenment for the architecture design of subsequent intelligent congestion control systems. Third, establishing an A/B test evaluation framework using a multi-dimensional QoE indicator system that can systematically and quantitatively compare the performance of different congestion control strategies under real network disturbances, providing a standardized evaluation methodology for the design of intelligent congestion control algorithms oriented to practical deployment.

Looking forward, there are several promising directions for future work. First, optimizing the reward function and policy bias to address the bandwidth underutilization issue in high-quality network scenarios. Second, expanding the scale and diversity of the training dataset to cover more device types, browser versions, and video content. Third, exploring multi-modal feature fusion by incorporating application-layer media features such as video frame quantization parameters and keyframe encoding sizes into the state space. Fourth, introducing federated learning mechanisms to enable continuous online fine-tuning of congestion control models without violating user privacy. Fifth, developing a more comprehensive QoE evaluation system that directly covers video subjective experience factors such as frame rate, resolution changes, and freezing duration.
}
